\section{Introdução}
\label{sec:introducao}

A proliferação de conteúdo multimídia, impulsionada pela onipresença da internet de alta velocidade e pela democratização das plataformas de criação, resultou em um volume de informação sem precedentes. Vídeos, podcasts e documentos digitais são produzidos e consumidos em uma escala exponencial, tornando-se fontes primárias de educação, entretenimento e comunicação profissional. No entanto, essa abundância de dados acarreta um desafio significativo: a sobrecarga de informação. Profissionais e estudantes frequentemente enfrentam a necessidade de assimilar o conteúdo de longas horas de material audiovisual ou de documentos extensos em um tempo limitado, um processo que é ineficiente e propenso à perda de foco.

O problema central abordado neste trabalho é a ineficiência no consumo de conteúdo de mídia de longa duração. A extração manual de informações essenciais de um vídeo de uma hora, por exemplo, exige que o espectador dedique uma hora de atenção contínua, o que nem sempre é viável ou produtivo. A ausência de ferramentas automatizadas para condensar esse conteúdo de forma inteligente representa uma barreira significativa para a acessibilidade e a produtividade na era da informação.

Para solucionar este problema, este trabalho propõe o design e a implementação de uma plataforma de software como serviço (SaaS) para o processamento e sumarização automática de mídias. A solução é uma API assíncrona capaz de receber diferentes tipos de mídia (vídeo, áudio, texto/PDF), orquestrar um pipeline de processamento em segundo plano e, ao final, entregar um resumo textual conciso do conteúdo original. O objetivo principal é reduzir drasticamente o tempo necessário para o consumo de informação, permitindo que o usuário absorva os pontos-chave de um material extenso em uma fração do tempo.

A relevância deste projeto reside não apenas na solução funcional para um problema contemporâneo, mas também na sua concepção arquitetural. Em contraste com abordagens monolíticas e rigidamente acopladas, a plataforma foi construída sobre uma arquitetura moderna, modular e orientada a eventos, empregando padrões de design avançados como CQRS (Command Query Responsibility Segregation), o Padrão State e um sistema de provedores dinâmicos. Esta abordagem garante não apenas a escalabilidade e a resiliência do sistema, mas também sua extensibilidade para suportar novos tipos de mídia e etapas de processamento no futuro, servindo como um estudo de caso prático para a construção de sistemas distribuídos complexos.

Este documento está estruturado da seguinte forma: a Seção \ref{sec:fundamentacao} apresenta a fundamentação teórica, detalhando os conceitos e tecnologias que sustentam o projeto. A Seção \ref{sec:arquitetura} (que você já tem) oferece uma visão aprofundada da arquitetura de software implementada. As seções subsequentes detalham o design, a implementação, os testes e a implantação da solução, culminando com a apresentação dos resultados e as conclusões do trabalho.